\mxnewpage\section*{EV4-T1 -- Préparation EAM}\label{EV4-T1}
\addcontentsline{toc}{section}{EV4-T1 -- Préparation EAM}
\textbf{Sujet d’entraînement}. Ensemble du programme.\par Durée : 2 h. Sans calculatrice. Note sur 20. Entraînements de mai et juin ; EAM de zone à confirmer.
\subsection*{Partie 1 -- Automatismes (6 points)}
Douze questions, \num{0,5} point chacune. Une seule réponse correcte ; aucune justification demandée. Une réponse incorrecte ou absente rapporte 0, sans retrait de points.
\begin{enumerate}
\item Calculer $5+2\times3$.\par A.~$14$\quad B.~$11$\quad C.~$21$\quad D.~$10$
\item Calculer $\frac12+\frac13$.\par A.~$\frac{5}{3}$\quad B.~$\frac{5}{6}$\quad C.~$\frac{1}{6}$\quad D.~$\frac{2}{5}$
\item Calculer $\dfrac{2^5}{2^{3}}$.\par A.~$1$\quad B.~$4$\quad C.~$10$\quad D.~$2$
\item Une valeur de 50 augmente de \pc{20}. Quelle est la valeur finale ?\par A.~$40$\quad B.~$10$\quad C.~$60$\quad D.~$70$
\item Une valeur de 80 est obtenue après une baisse de \pc{20}. Quelle est la valeur initiale ?\par A.~$60$\quad B.~$96$\quad C.~$100$\quad D.~$64$
\item Quel est le coefficient multiplicateur global d’une hausse de \pc{25} suivie d’une baisse de \pc{20} ?\par A.~$\frac{4}{5}$\quad B.~$1$\quad C.~$\frac{9}{10}$\quad D.~$\frac{21}{20}$
\item La solution de l’équation $2x+5=0$ est~:\par A.~$\frac{5}{2}$\quad B.~$10$\quad C.~$\frac{-5}{2}$\quad D.~$-10$
\item La pente de la droite passant par $\pt{0;5}$ et $\pt{2;11}$ est~:\par A.~$3$\quad B.~$-3$\quad C.~$6$\quad D.~$2$
\item La moyenne de $5, 7, 9$ est~:\par A.~$7$\quad B.~$21$\quad C.~$5$\quad D.~$8$
\item Un secteur de $90^\circ$ représente quelle fraction du total ?\par A.~$\frac{1}{4}$\quad B.~$\frac{3}{4}$\quad C.~$\frac{1}{3}$\quad D.~$\frac{1}{2}$
\item Une probabilité vaut $\dfrac{2}{5}$. Son complément est~:\par A.~$\frac{3}{5}$\quad B.~$\frac{1}{5}$\quad C.~$\frac{2}{5}$\quad D.~$\frac{5}{2}$
\item Quelle est la valeur de $f(-2)$ pour $f(x)=x^2+5$ ?\par A.~$20$\quad B.~$7$\quad C.~$1$\quad D.~$9$
\end{enumerate}
\mxnewpage\subsection*{Partie 2 -- Exercices indépendants (14 points)}
Toute démarche doit être justifiée. Les résultats exacts sont attendus sauf indication contraire.
\mxneed{12}
\subsection*{Exercice 1 -- Choisir une loi (5 points)}
On compare deux jeux. Le premier donne un gain $X=2$ avec certitude. Le second donne un gain $Y=-1$ avec probabilité $\dfrac23$ et un gain $Y=8$ avec probabilité $\dfrac13$. Les gains sont exprimés dans la même unité monétaire ; une valeur négative représente une perte.
\begin{enumerate}
\item Calculer $E(X)$ et $E(Y)$. \hfill(1 pt)
\item Calculer $V(Y)$ et comparer les dispersions des deux gains. \hfill(2 pts)
\item On joue aux deux jeux et l’on additionne les gains. Calculer $E(X+Y)$ en citant la propriété utilisée. \hfill(2 pts)
\end{enumerate}
\mxneed{12}
\subsection*{Exercice 2 -- Seuil discret et modèle continu (5 points)}
Une quantité est modélisée par $Q(t)=64e^{-t}$ pour $t\geq0$. On l’observe aux instants entiers $n\in\N$ et l’on pose $u_n=Q(n)$. Pour les comparaisons sans calculatrice, on donne l’encadrement $e^2<8<e^3$.
\begin{enumerate}
\item Étudier le sens de variation. \hfill(1 pt)
\item Déterminer le premier entier n tel que $Q(n)<8$ à l’aide de l’encadrement fourni. \hfill(2 pts)
\item Donner la nature de $u_n=Q(n)$ et l’expression de $u_0+\cdots+u_3$. \hfill(2 pts)
\end{enumerate}
\mxneed{12}
\subsection*{Exercice 3 -- Équation et cercle (4 points)}
Dans un repère orthonormé, on considère l’ensemble $\mathcal C$ défini par $x^2+y^2-2x-4y+1=0$. Il s’agit d’en déterminer les caractéristiques, puis d’identifier une tangente particulière.
\begin{enumerate}
\item Montrer que $\mathcal C$ est un cercle et préciser son centre et son rayon. \hfill(1 pt)
\item Déterminer ses points d’intersection avec la droite d’équation $x=1$. \hfill(1 pt)
\item Déterminer l’équation de la tangente au cercle au point de plus grande ordonnée et justifier. \hfill(2 pts)
\end{enumerate}
